\chapter{Cronología del Mundo de Ur-Bare}

\neofigure{transporter_bridge.jpg}{El histórico puente transbordador cruzando la Ría de Acero desde los primeros años de la colonia.}

La historia de NeoBilbao no es una crónica lineal de progreso tecnológico, sino el relato de una comunidad obrera que luchó por construir un hogar en los confines de la galaxia y acabó atrapada en los engranajes de un régimen corporativo monopolista. Esta es la cronología documentada de la colonia en el exoplaneta Ur-Bare.

\section{El éxodo y la Era del Auzolan (2052--2080)}

\begin{itemize}
  \item \lead{2052 -- Llegada de la flota Euskal-Astra:} Tras décadas de crisis ecológica y agotamiento de recursos en la Tierra, el consorcio cooperativo Euskal-Astra fleta una flota colonial que alcanza la órbita de \textbf{Ur-Bare}. El planeta es un páramo desértico, barrido por tempestades de polvo ferroso rojo y cubierto por una atmósfera tóxica irrespirable.
  
  \item \lead{2053 -- El hallazgo del Burdin-Urdin:} En las profundidades de una inmensa fosa tectónica bautizada como la \textbf{Ría de Acero}, los geólogos descubren concentraciones masivas de un mineral superconductor sin precedentes: el \emph{Burdin-Urdin} (Acero Azul). Su capacidad para canalizar densidades energéticas extremas sin resistencia lo convierte en el recurso más valioso de la galaxia para la navegación interestelar.
  
  \item \lead{2055--2075 -- La construcción comunitaria:} Los colonos desmontan el blindaje de las naves coloniales para encajarlas en los acantilados del cañón, fundando los primeros asentamientos protegidos. Para cruzar la sima sobre los ríos de metal fundido, ingenieros y soldadores levantan el colosal \textbf{Puente Transbordador}, inspirado en el de Portugalete. Durante estas tres décadas, la vida se organiza bajo el principio del \textbf{Auzolan} (trabajo vecinal solidario): no existen salarios ni castas, las raciones de aire y agua se comparten fraternalmente y las decisiones se toman en asamblea.
\end{itemize}

\section{El ascenso corporativo y La Margen (2080--2087)}

\begin{itemize}
  \item \lead{2080 -- La mutación de Orekak-Corp:} Originalmente contratada como empresa auxiliar de mantenimiento y logística, Orekak-Corp comienza a comprar licencias exclusivas de exportación de Burdin-Urdin y a monopolizar la fabricación de membranas de filtración de aire, debilitando el tejido cooperativo.
  
  \item \lead{2084 -- La consolidación del Distrito 7 («La Margen»):} Con la llegada masiva de nuevas oleadas de trabajadores, la ciudad se expande hacia los anillos medios de hormigón. Nace el sistema del \textbf{Índice Azul} y el publicitado \textbf{Programa Azul de Excelencia Ciudadana}, que promete un pase a la Cúpula Dorada para los ciudadanos más destacados. En las calles, los primeros agentes encubiertos (Los Auditores) comienzan a vigilar la disidencia política bajo un manto de aparente normalidad civil.
\end{itemize}

\section{La Guerra del Plasma y la privatización del aire (2088--2094)}

\neofigure{plasma_war_fighting.jpg}{Combates urbanos durante la Guerra del Plasma en los corredores industriales de la colonia.}

\begin{itemize}
  \item \lead{2088 -- El decreto del oxígeno tasado:} Orekak-Corp completa el domo inicial de contención atmosférica y proclama la privatización absoluta del aire filtrado. Respirar deja de ser un derecho colectivo y se convierte en una cuota mensual indexada a la producción de mineral.
  
  \item \lead{2089--2093 -- La Guerra de los Seis Años:} Las milicias obreras de la Ría responden declarando la huelga general y tomando las refinerías. Orekak-Corp militariza la seguridad corporativa: despliega armamento de plasma pesado, unidades anfibias en los colectores y los primeros exoesqueletos de minería pesada reconvertidos para el combate urbano (\textbf{Basajaun}).
  
  \item \lead{2094 -- La capitulación y la división de castas:} Asfixiados por el corte de suministros médicos y el racionamiento draconiano del oxígeno en las zonas rebeldes, los trabajadores se rinden. Orekak-Corp asume el poder dictatorial en NeoBilbao y sella formalmente la división entre la opulenta Cúpula Dorada y los Zócalos Industriales.
\end{itemize}

\section{La Gran Huelga y el Etxearen Zina (2102)}

\neofigure{workers_strike.jpg}{La histórica huelga obrera en la fundición de Burdin-Urdin antes del acuerdo del linaje.}

\begin{itemize}
  \item \lead{Marzo de 2102 -- La Huelga de las Dieciséis Horas:} Tras años de degradación laboral, jornadas interminables y reducción del caudal de aire en los niveles inferiores, los sindicatos clandestinos coordinan un apagón total de los altos hornos de Burdin-Urdin. Sin el flujo de energía de la fundición, la Cúpula Dorada se queda sin calefacción ni energía auxiliar.
  
  \item \lead{Abril de 2102 -- El pacto del Etxearen Zina:} Consciente de que no puede aniquilar a su propia mano de obra sin arruinarse, el Directorio corporativo formula una solución definitiva: ofrece salvar a la siguiente generación del veneno de las fundiciones. Los hijos de los trabajadores son reubicados en colegios e internados bioclimáticos de la Cúpula Dorada, con luz natural, aire puro y alimentación orgánica.
  
  \item \lead{El Juramento y el Velo:} A cambio de la salvación de sus hijos, los obreros (\emph{Langileak}) firman el \textbf{Etxearen Zina} (el Juramento del Linaje), comprometiéndose a turnos de servidumbre perpetua en los altos hornos. Para hacer soportable la separación, la corporación implanta a todos los trabajadores el \textbf{Velo}, una interfaz neuronal que les permite recibir transmisiones holográficas, cartas sensoriales y vídeos de sus hijos en cada descanso laboral, convirtiendo su sufrimiento en un sacrificio voluntario por el porvenir de su familia.
\end{itemize}
